% !TEX root = ../thesis.tex
% ==============================================================================
%  APPENDIX: 
%       TITLE PAGES & FORMS
%       (appendices/appendix_title_and_forms.tex)
% ==============================================================================
%  This appendix provides detailed architectural documentation on how the
%  Persian and English title pages, as well as official administrative forms
%  (signatures and declaration), are configured, formatted, and exported.
% ==============================================================================


\section{سازوکار صفحات عنوان و فرم‌های تصویب‌نامه و اظهارنامه}\label{sec:app_title_and_forms}

در معماری این قالب، تمامی اطلاعات شناسنامه‌ای، مشخصات سازمانی و اسامی هیئت داوران در یک فایل متمرکز به نشانی \code{frontmatter/info.tex} تعریف می‌شوند. این تفکیک ساختاری مانع از تکرار داده‌ها شده و تضمین می‌کند که هرگونه ویرایش در مشخصات پایان‌نامه، به‌صورت هماهنگ در چهار بخش سند اعمال گردد:
\begin{enumerate}
    \item صفحه‌ی عنوان فارسی روی جلد (\code{frontmatter/title.tex})؛
    \vspace{-0.25cm}
    \item صفحه‌ی تصویب‌نامه‌ی هیئت داوران (\code{frontmatter/signatures.tex})؛
    \vspace{-0.25cm}
    \item صفحه‌ی اظهارنامه‌ی اصالت متن (\code{frontmatter/declaration.tex})؛
    \vspace{-0.25cm}
    \item صفحه‌ی عنوان انگلیسی در انتهای سند (\code{backmatter/title\_en.tex}).
\end{enumerate}


\subsection{صفحات عنوان فارسی و انگلیسی و فواصل الاستیک}

صفحات عنوان بر اساس فنرهای فاصله‌ساز انعطاف‌پذیر\enfootnote{elastic fill glues} پیاده‌سازی شده‌اند. این سازوکار سبب می‌شود که فارغ از تک‌خطی یا چندخطی بودن عنوان پایان‌نامه یا تغییر در تعداد استادان، اجزای صفحه به‌طور خودکار و متوازن بر روی یک صفحه‌ی استاندارد \lr{A4} قرار گیرند. مقادیر پیش‌فرض این فواصل در فایل \code{styles/common.tex} تعریف شده و در صورت نیاز به سفارشی‌سازی فاصله‌ها، می‌توان مقادیر متناظر را در بخش ۳ فایل \code{styles/custom.tex} بازتعریف نمود.


\subsection{فرم‌های رسمی تصویب‌نامه و اظهارنامه}

فرم‌های اداری پایان‌نامه نیازمند رعایت الگوهای استاندارد دانشگاه صنعتی شریف هستند. در صفحه‌ی تصویب‌نامه، جدول اسامی داوران بر اساس شروط ماکروی \code{{\textbackslash}ifthesisfield} تنظیم می‌شود؛ به این معنا که در صورت عدم تعریف استاد مشاور یا داور دوم، ردیف‌های متناظر به‌طور خودکار حذف شده و قالب جدول دست‌نخورده باقی می‌ماند. همچنین ماکروی \code{{\textbackslash}ThesisSignTitle} این امکان را فراهم می‌سازد که اگر عنوان پایان‌نامه در صفحه‌ی روی جلد به دلیل شکست‌های دستی (\code{{\textbackslash}{\textbackslash}}) شکسته شده باشد، در فرم‌های اداری به‌صورت متنی پیوسته و روان درج گردد.


\subsection{تولید فرم‌های مستقل چاپی و ادغام نسخه‌های امضاشده}

برای تسهیل در آماده‌سازی مدارک، اسکریپت‌های ویژه‌ای در پوشه‌ی \code{scripts/} طراحی شده‌اند:
\begin{itemize}
    \item \textbf{تولید تک‌برگ فرم‌ها جهت چاپ}: با اجرای دستور زیر در ترمینال، فرم تصویب‌نامه به‌صورت یک فایل \lr{PDF} مستقل تک‌صفحه‌ای با سربرگ استاندارد تولید می‌شود:
    \begin{latin}
    \code{python scripts/export\_front\_pages.py signatures}
    \end{latin}
    \item \textbf{ادغام نسخه‌ی امضاشده پس از دفاع}: پس از برگزاری جلسه‌ی دفاع و اخذ امضای داوران، کافی است تصویر اسکن‌شده با فرمت \lr{PDF} در پوشه‌ی \code{frontmatter/} با نام \code{signatures.pdf} ذخیره شده و در فایل \code{thesis.tex} دستور \code{{\textbackslash}includepdfpages\{frontmatter/signatures.pdf\}} فعال گردد. این ماکرو فایل را با حفظ ابعاد دقیق \lr{A4} و بدون حاشیه جایگزین می‌نماید.
\end{itemize}

در جدول~\ref{tab:app_metadata_fields}، متغیرهای انتخابی تعبیه شده در فایل \code{frontmatter/info.tex} و کاربرد هر یک فهرست شده است.

\begin{table}[ht]
    \centering
    \scriptsize
    \setlength{\aboverulesep}{0pt}
    \setlength{\belowrulesep}{0pt}
    \caption{راهنمای متغیرهای فراداده در فایل \code{frontmatter/info.tex}}\label{tab:app_metadata_fields}
    \vspace{0.15cm}
    \begin{tabularx}{\textwidth}{l l l p{4cm}}
        \toprule
        \textbf{نام ماکرو} & \textbf{متغیر فارسی} & \textbf{معادل انگلیسی} & \textbf{شرح و کاربرد} \\
        \midrule
        \code{{\textbackslash}ThesisDegree} & \code{{\textbackslash}ThesisDegree} & \code{{\textbackslash}EnglishThesisDegree} & مقطع (کارشناسی ارشد، دکتری) \\
        \code{{\textbackslash}ThesisType} & \code{{\textbackslash}ThesisType} & \code{{\textbackslash}EnglishThesisType} & نوع (پایان‌نامه، رساله، گزارش) \\
        \code{{\textbackslash}ThesisUniversity} & \code{{\textbackslash}ThesisUniversity} & \code{{\textbackslash}EnglishThesisUniversity} & نام رسمی دانشگاه \\
        \code{{\textbackslash}ThesisDepartment} & \code{{\textbackslash}ThesisDepartment} & \code{{\textbackslash}EnglishThesisDepartment} & نام دانشکده‌ی محل تحصیل \\
        \code{{\textbackslash}ThesisMajor} & \code{{\textbackslash}ThesisMajor} & \code{{\textbackslash}EnglishThesisMajor} & رشته و گرایش تحصیلی \\
        \code{{\textbackslash}ThesisTitle} & \code{{\textbackslash}ThesisTitle} & \code{{\textbackslash}EnglishThesisTitle} & عنوان پایان‌نامه روی جلد \\
        \code{{\textbackslash}ThesisAuthor} & \code{{\textbackslash}ThesisAuthor} & \code{{\textbackslash}EnglishThesisAuthor} & نام و نام خانوادگی دانشجو \\
        \code{{\textbackslash}ThesisSupervisor} & \code{{\textbackslash}ThesisSupervisor} & \code{{\textbackslash}EnglishThesisSupervisor} & نام و مرتبه‌ی علمی استاد راهنما \\
        \code{{\textbackslash}ThesisAdvisor} & \code{{\textbackslash}ThesisAdvisor} & \code{{\textbackslash}EnglishThesisAdvisor} & نام استاد مشاور اول (اختیاری) \\
        \code{{\textbackslash}ThesisSecondAdvisor} & \code{{\textbackslash}ThesisSecondAdvisor} & \code{{\textbackslash}EnglishThesisSecondAdvisor} & نام استاد مشاور دوم (اختیاری) \\
        \code{{\textbackslash}ThesisExaminer} & \code{{\textbackslash}ThesisExaminer} & \code{{\textbackslash}EnglishThesisExaminer} & نام استاد ممتحن اول (اختیاری) \\
        \code{{\textbackslash}ThesisSecondExaminer} & \code{{\textbackslash}ThesisSecondExaminer} & \code{{\textbackslash}EnglishThesisSecondExaminer} & نام استاد ممتحن دوم (اختیاری) \\
        \code{{\textbackslash}ThesisThirdExaminer} & \code{{\textbackslash}ThesisThirdExaminer} & \code{{\textbackslash}EnglishThesisThirdExaminer} & نام استاد ممتحن سوم (اختیاری) \\
        \code{{\textbackslash}ThesisDate} & \code{{\textbackslash}ThesisDate} & \code{{\textbackslash}EnglishThesisDate} & تاریخ دفاع برای صفحات عنوان \\
        \code{{\textbackslash}ThesisSignDate} & \code{{\textbackslash}ThesisSignDate} & --- & تاریخ امضای فرم‌های رسمی \\
        \code{{\textbackslash}ThesisSignTitle} & \code{{\textbackslash}ThesisSignTitle} & --- & عنوان پیوسته برای فرم‌های اداری \\
        \bottomrule
    \end{tabularx}\vspace{0.5cm}
\end{table}
